% --- IMPORTS ---
\documentclass[leqno]{article}
\usepackage{graphicx}
\usepackage{dsfont}
\usepackage{mathtools}
\usepackage{amssymb}
\usepackage{amsmath}
\usepackage{amsthm}
\usepackage{float}
\usepackage{ragged2e}
\usepackage{tensor}
\usepackage{fancyhdr}
\usepackage{empheq}
\usepackage[most]{tcolorbox}
\usepackage{enumitem}
\usepackage{dirtytalk}
\usepackage{marginnote}
\usepackage{microtype}
\usepackage[
  a4paper,
  left=0.8in,
  right=1in,
  top=1in,
  bottom=0.5in,
  headheight=14pt,
  headsep=0.4in
]{geometry}
\setlength{\parindent}{0cm}
% --- TikZ Packages ---
\usepackage{pgfplots}
\pgfplotsset{compat=1.18}
\usepgfplotslibrary{fillbetween, polar}
\usetikzlibrary{calc, positioning}
\usetikzlibrary{angles, quotes}
\usetikzlibrary{arrows.meta}
\usetikzlibrary{decorations.pathreplacing, decorations.markings}
\usetikzlibrary{patterns}
\usetikzlibrary{intersections}
% --- URL Packages ---
\usepackage[colorlinks=true,linkcolor=red,citecolor=magenta,urlcolor=black]{hyperref}%
\urlstyle{same}
% --- END OF IMPORTS ---

% --- CUSTOM COMMANDS ---
\RenewDocumentCommand{\marks}{o m}{%
  \IfValueTF{#1}%
    {\marginnote{\raggedright\textbf{[#2]}}[#1]}%
    {\marginnote{\raggedright\textbf{[#2]}}}%
}
\newcommand{\vspacerule}[2]{%
  \par\vspace*{#1}%
  \noindent\hrulefill\par
  \vspace*{#2}%
}
\definecolor{scarlet}{RGB}{205, 40, 75}
\newcommand{\questionitem}{%
  \item\phantomsection
  \addcontentsline{toc}{section}{Question \number\value{enumi}}%
}
\renewcommand{\contentsname}{Questions}
% --- END OF CUSTOM COMMANDS ---

% --- HEADER CONFIG ---
\pagestyle{fancy}
\fancyhead{}
\fancyfoot{}
\renewcommand{\headrulewidth}{0.9pt}
\lhead{\textit{Solutions} - Roots of Polynomials}
\chead{}
\rhead{\href{https://danieljsmith.org}{danieljsmith.org}}
% --- END OF HEADER CONFIG ---

\begin{document}
\vspace*{20pt}
\tableofcontents
\newpage
\begin{enumerate}[
  label=\textbf{\arabic*.},
  leftmargin=*,
  topsep=0pt,
  partopsep=0pt,
  parsep=0pt
]

% ---- Question 1 ----
\questionitem
% [Source: _QBT___Solns__Roots_of_Polynomials, Q1]

The cubic equation
\[
x^3 - 5x^2 + px - q = 0
\]
where $p$ and $q$ are positive real constants, has roots $\alpha$, $\beta$ and $\gamma$. \\

Given that
\[
(\alpha-\beta)^2 + (\beta-\gamma)^2 + (\gamma-\alpha)^2 = 14
\]
\begin{enumerate}[label=\textbf{(\alph*)}]
  \item show that $p=6$ \marks{3} \\
\end{enumerate}
Given that
\[
\frac{1}{\alpha} + \frac{1}{\beta} + \frac{1}{\gamma} = 3
\]
\begin{enumerate}[label=\textbf{(\alph*)}, resume]
  \item determine the value of $q$ \marks{3} \\
  \item Without solving the cubic equation, determine the value of $(\alpha+1)(\beta+1)(\gamma+1)$ \marks{4} \\
\end{enumerate}

\vspace*{10pt}

\begin{tcolorbox}[colback=white, colframe=scarlet, title={\textbf{Solution}}, breakable, grow to left by=6mm, grow to right by=10mm]
\begin{enumerate}[label=\textbf{(\alph*)}]
  \item For the cubic
  \[
  x^3-5x^2+px-q=0
  \]
  with roots $\alpha,\beta,\gamma$, the relations between roots and coefficients give
  \[
  \alpha+\beta+\gamma=5,
  \qquad
  \alpha\beta+\beta\gamma+\gamma\alpha=p
  \]

  Now expand the given expression:
  \begin{align*}
  (\alpha-\beta)^2+(\beta-\gamma)^2+(\gamma-\alpha)^2
  &= (\alpha^2+\beta^2-2\alpha\beta)+(\beta^2+\gamma^2-2\beta\gamma)+(\gamma^2+\alpha^2-2\gamma\alpha) \\
  &= 2(\alpha^2+\beta^2+\gamma^2)-2(\alpha\beta+\beta\gamma+\gamma\alpha)
  \end{align*}

  Also,
  \[
  \alpha^2+\beta^2+\gamma^2=(\alpha+\beta+\gamma)^2-2(\alpha\beta+\beta\gamma+\gamma\alpha)
  \]

  So
  \begin{align*}
  (\alpha-\beta)^2+(\beta-\gamma)^2+(\gamma-\alpha)^2
  &=2\big((\alpha+\beta+\gamma)^2-2(\alpha\beta+\beta\gamma+\gamma\alpha)\big)-2(\alpha\beta+\beta\gamma+\gamma\alpha) \\
  &=2\big((\alpha+\beta+\gamma)^2-3(\alpha\beta+\beta\gamma+\gamma\alpha)\big)
  \end{align*}

  Substituting the known values,
  \begin{align*}
  14&=2(5^2-3p) \\
  14&=2(25-3p) \\
  7&=25-3p \\
  3p&=18 \\
  p&=6
  \end{align*}

  Hence, $p=6$.

  \item For this cubic,
  \[
  \alpha\beta\gamma=q
  \]

  Also,
  \begin{align*}
  \frac{1}{\alpha}+\frac{1}{\beta}+\frac{1}{\gamma}
  &=\frac{\beta\gamma+\gamma\alpha+\alpha\beta}{\alpha\beta\gamma} \\
  &=\frac{p}{q}
  \end{align*}

  Given that this sum is $3$,
  \[
  \frac{p}{q}=3
  \]

  From part (a), $p=6$, so
  \begin{align*}
  \frac{6}{q}&=3 \\
  6&=3q \\
  q&=2
  \end{align*}

  Hence, $q=2$.

  \item Expand:
  \begin{align*}
  (\alpha+1)(\beta+1)(\gamma+1)
  &=\alpha\beta\gamma+(\alpha\beta+\beta\gamma+\gamma\alpha)+(\alpha+\beta+\gamma)+1
  \end{align*}

  Using the root relations,
  \[
  \alpha\beta\gamma=q,\qquad
  \alpha\beta+\beta\gamma+\gamma\alpha=p,\qquad
  \alpha+\beta+\gamma=5
  \]

  Therefore
  \begin{align*}
  (\alpha+1)(\beta+1)(\gamma+1)
  &=q+p+5+1 \\
  &=2+6+5+1 \\
  &=14
  \end{align*}

  Hence,
  \[
  (\alpha+1)(\beta+1)(\gamma+1)=14
  \]
\end{enumerate}
\end{tcolorbox}


\newpage

% ---- Question 2 ----
\questionitem
% [Source: _QBT___Solns__Roots_of_Polynomials, Q2]

The quartic equation $x^4+bx^3+cx^2+dx+3=0$ has roots $\alpha$, $\beta$, $\gamma$, $\delta$. It is given that
\[
\alpha+\beta+\gamma+\delta=-2,\qquad (\alpha+1)^2+(\beta+1)^2+(\gamma+1)^2+(\delta+1)^2=6,\qquad \alpha^{-1}+\beta^{-1}+\gamma^{-1}+\delta^{-1}=2
\]
\begin{enumerate}[label=\textbf{(\alph*)}]
  \item Find the values of $b$, $c$ and $d$ \marks{6} \\
  \item Given also that
  \[
  (\alpha+1)^3+(\beta+1)^3+(\gamma+1)^3+(\delta+1)^3=20
  \]
  find the value of $\alpha^4+\beta^4+\gamma^4+\delta^4$ \marks{2} \\
\end{enumerate}

\vspace*{10pt}

\begin{tcolorbox}[colback=white, colframe=scarlet, title={\textbf{Solution}}, breakable, grow to left by=6mm, grow to right by=10mm]
\begin{enumerate}[label=\textbf{(\alph*)}]
  \item For the monic quartic
  \[
  x^4+bx^3+cx^2+dx+3=0
  \]
  with roots $\alpha,\beta,\gamma,\delta$, the standard relations between coefficients and roots give
  \[
  \alpha+\beta+\gamma+\delta=-b
  \]
  \[
  \alpha\beta+\alpha\gamma+\alpha\delta+\beta\gamma+\beta\delta+\gamma\delta=c
  \]
  \[
  \alpha\beta\gamma+\alpha\beta\delta+\alpha\gamma\delta+\beta\gamma\delta=-d
  \]
  \[
  \alpha\beta\gamma\delta=3
  \]

  We are given
  \[
  \alpha+\beta+\gamma+\delta=-2
  \]
  so
  \[
  -b=-2
  \]
  hence
  \[
  b=2
  \]

  Next,
  \[
  \alpha^{-1}+\beta^{-1}+\gamma^{-1}+\delta^{-1}
  =\frac{\beta\gamma\delta+\alpha\gamma\delta+\alpha\beta\delta+\alpha\beta\gamma}{\alpha\beta\gamma\delta}
  \]
  Therefore
  \[
  2=\frac{-d}{3}
  \]
  so
  \[
  d=-6
  \]

  Now use
  \[
  (\alpha+1)^2+(\beta+1)^2+(\gamma+1)^2+(\delta+1)^2=6
  \]
  Expanding,
  \[
  \alpha^2+\beta^2+\gamma^2+\delta^2+2(\alpha+\beta+\gamma+\delta)+4=6
  \]
  since there are $4$ lots of $1^2$.

  Substituting $\alpha+\beta+\gamma+\delta=-2$,
  \[
  \alpha^2+\beta^2+\gamma^2+\delta^2+2(-2)+4=6
  \]
  \[
  \alpha^2+\beta^2+\gamma^2+\delta^2=6
  \]

  Also,
  \[
  (\alpha+\beta+\gamma+\delta)^2
  =\alpha^2+\beta^2+\gamma^2+\delta^2+2(\alpha\beta+\alpha\gamma+\alpha\delta+\beta\gamma+\beta\delta+\gamma\delta)
  \]
  so
  \[
  (-2)^2=6+2c
  \]
  \[
  4=6+2c
  \]
  \[
  2c=-2
  \]
  \[
  c=-1
  \]

  Therefore,
  \[
  b=2,\qquad c=-1,\qquad d=-6
  \]

  \item Using the extra condition
  \[
  (\alpha+1)^3+(\beta+1)^3+(\gamma+1)^3+(\delta+1)^3=20
  \]
  expand:
  \[
  \sum (\text{root}+1)^3=\sum \text{root}^3+3\sum \text{root}^2+3\sum \text{root}+4
  \]
  Hence
  \[
  \alpha^3+\beta^3+\gamma^3+\delta^3+3(\alpha^2+\beta^2+\gamma^2+\delta^2)+3(\alpha+\beta+\gamma+\delta)+4=20
  \]
  Using $\alpha^2+\beta^2+\gamma^2+\delta^2=6$ and $\alpha+\beta+\gamma+\delta=-2$,
  \[
  \alpha^3+\beta^3+\gamma^3+\delta^3+3(6)+3(-2)+4=20
  \]
  \[
  \alpha^3+\beta^3+\gamma^3+\delta^3+18-6+4=20
  \]
  \[
  \alpha^3+\beta^3+\gamma^3+\delta^3=4
  \]

  From part (a), the quartic is
  \[
  x^4+2x^3-x^2-6x+3=0
  \]
  so each root $r$ satisfies
  \[
  r^4+2r^3-r^2-6r+3=0
  \]
  and therefore
  \[
  r^4=-2r^3+r^2+6r-3
  \]

  Summing this for $r=\alpha,\beta,\gamma,\delta$ gives
  \[
  \alpha^4+\beta^4+\gamma^4+\delta^4
  =-2(\alpha^3+\beta^3+\gamma^3+\delta^3)+(\alpha^2+\beta^2+\gamma^2+\delta^2)+6(\alpha+\beta+\gamma+\delta)-12
  \]
  \[
  =-2(4)+6+6(-2)-12
  \]
  \[
  =-8+6-12-12
  \]
  \[
  =-26
  \]

  Therefore,
  \[
  \alpha^4+\beta^4+\gamma^4+\delta^4=-26
  \]
\end{enumerate}
\end{tcolorbox}


\newpage

% ---- Question 3 ----
\questionitem
% [Source: _QBT___Solns__Roots_of_Polynomials, Q3]

The equation $x^3 - 4x^2 - x + 6 = 0$ has roots $\alpha$, $\beta$ and $\gamma$. \\

Determine the value of \[\frac{1}{\alpha+\beta} + \frac{1}{\beta+\gamma} + \frac{1}{\gamma+\alpha}\] \marks{4}

\vspace*{10pt}

\begin{tcolorbox}[colback=white, colframe=scarlet, title={\textbf{Solution}}, breakable, grow to left by=6mm, grow to right by=10mm]
Let
\[
s_1=\alpha+\beta+\gamma,\qquad s_2=\alpha\beta+\beta\gamma+\gamma\alpha,\qquad s_3=\alpha\beta\gamma
\]

From
\[
x^3-4x^2-x+6=0
\]
comparing with \(x^3-s_1x^2+s_2x-s_3=0\), we get
\[
s_1=4,\qquad s_2=-1,\qquad s_3=-6
\]

We need
\[
\frac{1}{\alpha+\beta}+\frac{1}{\beta+\gamma}+\frac{1}{\gamma+\alpha}
\]

Taking a common denominator,
\begin{align*}
\frac{1}{\alpha+\beta}+\frac{1}{\beta+\gamma}+\frac{1}{\gamma+\alpha}
&=\frac{(\beta+\gamma)+(\gamma+\alpha)+(\alpha+\beta)}
{(\alpha+\beta)(\beta+\gamma)(\gamma+\alpha)}\\
&=\frac{2(\alpha+\beta+\gamma)}
{(\alpha+\beta)(\beta+\gamma)(\gamma+\alpha)}\\
&=\frac{2s_1}{(\alpha+\beta)(\beta+\gamma)(\gamma+\alpha)}
\end{align*}

Now expand the denominator:
\begin{align*}
(\alpha+\beta)(\beta+\gamma)(\gamma+\alpha)
&=(\alpha+\beta)\bigl(\beta\gamma+\beta\alpha+\gamma^2+\gamma\alpha\bigr)\\
&=\alpha^2\beta+\alpha^2\gamma+\alpha\beta^2+\beta^2\gamma+\alpha\gamma^2+\beta\gamma^2+2\alpha\beta\gamma
\end{align*}

Also,
\begin{align*}
(\alpha+\beta+\gamma)(\alpha\beta+\beta\gamma+\gamma\alpha)
&=\alpha^2\beta+\alpha^2\gamma+\alpha\beta^2+\beta^2\gamma+\alpha\gamma^2+\beta\gamma^2+3\alpha\beta\gamma
\end{align*}

So
\[
(\alpha+\beta)(\beta+\gamma)(\gamma+\alpha)=s_1s_2-s_3
\]

Hence
\[
\frac{1}{\alpha+\beta}+\frac{1}{\beta+\gamma}+\frac{1}{\gamma+\alpha}
=\frac{2s_1}{s_1s_2-s_3}
\]

Substitute \(s_1=4\), \(s_2=-1\), \(s_3=-6\):
\begin{align*}
\frac{2s_1}{s_1s_2-s_3}
&=\frac{2(4)}{4(-1)-(-6)}\\
&=\frac{8}{-4+6}\\
&=\frac{8}{2}\\
&=4
\end{align*}

Therefore, the required value is \(4\).
\end{tcolorbox}


\newpage

% ---- Question 4 ----
\questionitem
% [Source: _QBT___Solns__Roots_of_Polynomials, Q4]

The quartic equation $3x^4 - x^3 - 10x^2 + 8x + 4 = 0$ has roots $\alpha$, $\beta$, $\gamma$ and $\delta$. \\

Without solving the equation, find equations with integer coefficients whose roots are \\
\begin{enumerate}[label=\textbf{(\alph*)}]
  \item $\dfrac{1}{\alpha}$, $\dfrac{1}{\beta}$, $\dfrac{1}{\gamma}$ and $\dfrac{1}{\delta}$ \marks{6} \\
  \item $\dfrac{2}{\alpha - 1}$, $\dfrac{2}{\beta - 1}$, $\dfrac{2}{\gamma - 1}$ and $\dfrac{2}{\delta - 1}$ \marks{6} \\
\end{enumerate}

\vspace*{10pt}

\begin{tcolorbox}[colback=white, colframe=scarlet, title={\textbf{Solution}}, breakable, grow to left by=6mm, grow to right by=10mm]
\begin{enumerate}[label=\textbf{(\alph*)}]
  \item Let
\[
f(x)=3x^4-x^3-10x^2+8x+4
\]
and let \(x=\alpha\) be any root, so \(f(\alpha)=0\).

We want an equation whose roots are
\[
\frac1\alpha,\ \frac1\beta,\ \frac1\gamma,\ \frac1\delta
\]

Since the constant term of \(f(x)\) is \(4\), none of the roots is \(0\), so reciprocals are defined.

Let
\[
y=\frac1x
\]
so that
\[
x=\frac1y
\]

Substitute \(x=\dfrac1y\) into \(f(x)=0\):
\[
3\left(\frac1y\right)^4-\left(\frac1y\right)^3-10\left(\frac1y\right)^2+8\left(\frac1y\right)+4=0
\]

This gives
\[
\frac{3}{y^4}-\frac{1}{y^3}-\frac{10}{y^2}+\frac{8}{y}+4=0
\]

Multiply through by \(y^4\):
\[
3-y-10y^2+8y^3+4y^4=0
\]

Rearrange in descending powers:
\[
4y^4+8y^3-10y^2-y+3=0
\]

So the equation whose roots are
\(\dfrac1\alpha,\dfrac1\beta,\dfrac1\gamma,\dfrac1\delta\) is
\[
4x^4+8x^3-10x^2-x+3=0
\]

  \item We now want an equation whose roots are
\[
\frac{2}{\alpha-1},\ \frac{2}{\beta-1},\ \frac{2}{\gamma-1},\ \frac{2}{\delta-1}
\]

First check that \(1\) is not a root of the original equation:
\[
f(1)=3-1-10+8+4=4\neq 0
\]
so \(\alpha-1,\beta-1,\gamma-1,\delta-1\) are all non-zero.

Let
\[
y=\frac{2}{x-1}
\]
Then
\[
x-1=\frac{2}{y}
\qquad\Rightarrow\qquad
x=1+\frac{2}{y}
\]

Substitute \(x=1+\dfrac{2}{y}\) into \(f(x)=0\):
\[
3\left(1+\frac{2}{y}\right)^4-\left(1+\frac{2}{y}\right)^3-10\left(1+\frac{2}{y}\right)^2+8\left(1+\frac{2}{y}\right)+4=0
\]

Now expand each power:
\begin{align*}
\left(1+\frac{2}{y}\right)^2&=1+\frac{4}{y}+\frac{4}{y^2} \\
\left(1+\frac{2}{y}\right)^3&=1+\frac{6}{y}+\frac{12}{y^2}+\frac{8}{y^3} \\
\left(1+\frac{2}{y}\right)^4&=1+\frac{8}{y}+\frac{24}{y^2}+\frac{32}{y^3}+\frac{16}{y^4}
\end{align*}

So
\begin{align*}
0={}&3\left(1+\frac{8}{y}+\frac{24}{y^2}+\frac{32}{y^3}+\frac{16}{y^4}\right)
-\left(1+\frac{6}{y}+\frac{12}{y^2}+\frac{8}{y^3}\right) \\
&-10\left(1+\frac{4}{y}+\frac{4}{y^2}\right)
+8\left(1+\frac{2}{y}\right)+4
\end{align*}

Expand fully:
\begin{align*}
0={}&\left(3+\frac{24}{y}+\frac{72}{y^2}+\frac{96}{y^3}+\frac{48}{y^4}\right)
-\left(1+\frac{6}{y}+\frac{12}{y^2}+\frac{8}{y^3}\right) \\
&-\left(10+\frac{40}{y}+\frac{40}{y^2}\right)
+\left(8+\frac{16}{y}\right)+4
\end{align*}

Collect like terms:
\begin{align*}
0&=\left(3-1-10+8+4\right)
+\left(\frac{24}{y}-\frac{6}{y}-\frac{40}{y}+\frac{16}{y}\right) \\
&\quad+\left(\frac{72}{y^2}-\frac{12}{y^2}-\frac{40}{y^2}\right)
+\left(\frac{96}{y^3}-\frac{8}{y^3}\right)
+\frac{48}{y^4}
\end{align*}

Hence
\[
0=4-\frac{6}{y}+ \frac{20}{y^2}+\frac{88}{y^3}+\frac{48}{y^4}
\]

Multiply through by \(y^4\):
\[
4y^4-6y^3+20y^2+88y+48=0
\]

All coefficients are even, so divide by \(2\):
\[
2y^4-3y^3+10y^2+44y+24=0
\]

Therefore the equation whose roots are
\(\dfrac{2}{\alpha-1},\dfrac{2}{\beta-1},\dfrac{2}{\gamma-1},\dfrac{2}{\delta-1}\) is
\[
2x^4-3x^3+10x^2+44x+24=0
\]
\end{enumerate}
\end{tcolorbox}


\newpage

% ---- Question 5 ----
\questionitem
% [Source: _QBT___Solns__Roots_of_Polynomials, Q5]

The cubic equation $x^3 + x^2 - 2x - 1 = 0$ has roots $\alpha$, $\beta$ and $\gamma$. \\

\begin{enumerate}[label=\textbf{(\alph*)}]
  \item Find a cubic equation whose roots are $\frac{1}{\alpha+1}$, $\frac{1}{\beta+1}$ and $\frac{1}{\gamma+1}$ \marks{3} \\
  \item Find the value of \[\left(\frac{1}{\alpha+1}\right)^2 + \left(\frac{1}{\beta+1}\right)^2 + \left(\frac{1}{\gamma+1}\right)^2\] \marks[-20pt]{2} \\
  \item Find the value of \[\left(\frac{1}{\alpha+1}\right)^3 + \left(\frac{1}{\beta+1}\right)^3 + \left(\frac{1}{\gamma+1}\right)^3\] \marks{2} \\
\end{enumerate}

\vspace*{10pt}

\begin{tcolorbox}[colback=white, colframe=scarlet, title={\textbf{Solution}}, breakable, grow to left by=6mm, grow to right by=10mm]
\begin{enumerate}[label=\textbf{(\alph*)}]
  \item Since $x=-1$ is not a root of $x^3+x^2-2x-1=0$, the quantities
  \[
  \frac{1}{\alpha+1},\quad \frac{1}{\beta+1},\quad \frac{1}{\gamma+1}
  \]
  are all defined.

  Let
  \[
  y=\frac{1}{x+1}
  \]
  so
  \[
  x=\frac{1}{y}-1=\frac{1-y}{y}
  \]

  Substitute into the given cubic:
  \begin{align*}
  \left(\frac{1-y}{y}\right)^3+\left(\frac{1-y}{y}\right)^2-2\left(\frac{1-y}{y}\right)-1&=0
  \end{align*}

  Multiply through by $y^3$:
  \begin{align*}
  (1-y)^3+y(1-y)^2-2y^2(1-y)-y^3&=0
  \end{align*}

  Expanding:
  \begin{align*}
  (1-3y+3y^2-y^3)+(y-2y^2+y^3)-2y^2+2y^3-y^3&=0
  \end{align*}

  Collecting terms:
  \begin{align*}
  1-2y-y^2+y^3&=0
  \end{align*}

  So
  \[
  y^3-y^2-2y+1=0
  \]

  When $x=\alpha,\beta,\gamma$, the corresponding values of $y$ are
  \[
  \frac{1}{\alpha+1},\quad \frac{1}{\beta+1},\quad \frac{1}{\gamma+1}
  \]
  so the required cubic equation is
  \[
  y^3-y^2-2y+1=0
  \]

  \item Let the roots of
  \[
  y^3-y^2-2y+1=0
  \]
  be $p,q,r$, where
  \[
  p=\frac{1}{\alpha+1},\quad q=\frac{1}{\beta+1},\quad r=\frac{1}{\gamma+1}
  \]

  From the cubic,
  \[
  p+q+r=1,\qquad pq+pr+qr=-2
  \]

  Now
  \begin{align*}
  p^2+q^2+r^2&=(p+q+r)^2-2(pq+pr+qr)\\
  &=1^2-2(-2)\\
  &=5
  \end{align*}

  Therefore
  \[
  \left(\frac{1}{\alpha+1}\right)^2+\left(\frac{1}{\beta+1}\right)^2+\left(\frac{1}{\gamma+1}\right)^2=5
  \]

  \item Using the same roots $p,q,r$, we also have
  \[
  pqr=-1
  \]

  Use
  \[
  p^3+q^3+r^3=(p+q+r)^3-3(p+q+r)(pq+pr+qr)+3pqr
  \]

  Hence
  \begin{align*}
  p^3+q^3+r^3&=1^3-3(1)(-2)+3(-1)\\
  &=1+6-3\\
  &=4
  \end{align*}

  Therefore
  \[
  \left(\frac{1}{\alpha+1}\right)^3+\left(\frac{1}{\beta+1}\right)^3+\left(\frac{1}{\gamma+1}\right)^3=4
  \]
\end{enumerate}
\end{tcolorbox}


\newpage

% ---- Question 6 ----
\questionitem
% [Source: _QBT___Solns__Roots_of_Polynomials, Q6]

The equation $x^3 - 14x^2 + 56x + c = 0$, where $c$ is a constant, has three roots in geometric progression. \\

\begin{enumerate}[label=\textbf{(\alph*)}]
  \item Determine the roots of the equation \marks{6} \\
  \item Find $c$ \marks{1} \\
\end{enumerate}

\vspace*{10pt}

\begin{tcolorbox}[colback=white, colframe=scarlet, title={\textbf{Solution}}, breakable, grow to left by=6mm, grow to right by=10mm]
\begin{enumerate}[label=\textbf{(\alph*)}]
\item Let the three roots be
\[
\frac{a}{r},\quad a,\quad ar
\]
where $r\neq 0$.

For the cubic
\[
x^3-14x^2+56x+c=0
\]
the sum of the roots is $14$, and the sum of the products of the roots taken two at a time is $56$.

So
\begin{align*}
\frac{a}{r}+a+ar&=14 \\
a\left(\frac1r+1+r\right)&=14
\end{align*}

Also,
\begin{align*}
\left(\frac{a}{r}\right)(a)+a(ar)+\left(\frac{a}{r}\right)(ar)&=56 \\
\frac{a^2}{r}+a^2r+a^2&=56 \\
a^2\left(\frac1r+1+r\right)&=56
\end{align*}

Now divide the second equation by the first:
\[
\frac{a^2\left(\frac1r+1+r\right)}{a\left(\frac1r+1+r\right)}=\frac{56}{14}
\]
so
\[
a=4
\]

Substitute into
\[
a\left(\frac1r+1+r\right)=14
\]
to get
\begin{align*}
4\left(\frac1r+1+r\right)&=14 \\
\frac1r+1+r&=\frac72
\end{align*}

Multiply through by $2r$:
\begin{align*}
2+2r+2r^2&=7r \\
2r^2-5r+2&=0
\end{align*}

Factorising,
\[
(2r-1)(r-2)=0
\]
so
\[
r=\frac12 \quad \text{or} \quad r=2
\]

Hence the roots are
\[
\frac{4}{2},\ 4,\ 4(2)
\]
that is
\[
2,\ 4,\ 8
\]

So the roots are $\boxed{2,\,4,\,8}$.

\item The product of the roots is
\[
2\cdot 4\cdot 8=64
\]

For
\[
x^3-14x^2+56x+c=0
\]
the product of the roots is $-c$.

So
\[
-c=64
\]
and therefore
\[
c=-64
\]

So $\boxed{c=-64}$.
\end{enumerate}
\end{tcolorbox}


\newpage

% ---- Question 7 ----
\questionitem
% [Source: _QBT___Solns__Roots_of_Polynomials, Q7]

The quadratic equation
\[
x^2 - kx + 1 = 0
\]
where $k$ is an integer, has roots $\alpha$ and $\beta$. \\

\begin{enumerate}[label=\textbf{(\alph*)}]
  \item Write down, in terms of $k$ where appropriate, the value of $\alpha + \beta$ and the value of $\alpha\beta$ \marks{2} \\
  \item Determine, in simplest form in terms of $k$, the value of
  \[
  (\alpha^2 + \beta) + (\beta^2 + \alpha)
  \]
  \marks[-20pt]{3}
  \item Determine a quadratic equation which has roots
  \[
  \alpha^2 + \beta \text{ and } \beta^2 + \alpha
  \]
  giving your answer in the form $px^2 + qx + r = 0$, where $p$, $q$ and $r$ are integer values in terms of $k$ \marks{4} \\
\end{enumerate}

\vspace*{10pt}

\begin{tcolorbox}[colback=white, colframe=scarlet, title={\textbf{Solution}}, breakable, grow to left by=6mm, grow to right by=10mm]
\begin{enumerate}[label=\textbf{(\alph*)}]
  \item For the quadratic \(x^2-kx+1=0\), we use
  \[
  \alpha+\beta=-\frac{\text{coefficient of }x}{\text{coefficient of }x^2},
  \qquad
  \alpha\beta=\frac{\text{constant term}}{\text{coefficient of }x^2}
  \]
  So
  \[
  \alpha+\beta=-\frac{-k}{1}=k,
  \qquad
  \alpha\beta=\frac{1}{1}=1
  \]
  Hence
  \[
  \alpha+\beta=k,\qquad \alpha\beta=1
  \]

  \item
  \begin{align*}
  (\alpha^2+\beta)+(\beta^2+\alpha)
  &=\alpha^2+\beta^2+\alpha+\beta\\
  &=(\alpha+\beta)^2-2\alpha\beta+(\alpha+\beta)\\
  &=k^2-2(1)+k\\
  &=k^2+k-2
  \end{align*}
  Hence the value is
  \[
  k^2+k-2
  \]

  \item Let the new roots be
  \[
  \alpha^2+\beta \quad \text{and} \quad \beta^2+\alpha
  \]
  Their sum is, from part (b),
  \[
  (\alpha^2+\beta)+(\beta^2+\alpha)=k^2+k-2
  \]

  Now find their product:
  \begin{align*}
  (\alpha^2+\beta)(\beta^2+\alpha)
  &=\alpha^2\beta^2+\alpha^3+\beta^3+\alpha\beta\\
  &=(\alpha\beta)^2+\alpha^3+\beta^3+\alpha\beta\\
  &=1+\alpha^3+\beta^3+1\\
  &=\alpha^3+\beta^3+2
  \end{align*}

  Next,
  \begin{align*}
  \alpha^3+\beta^3
  &=(\alpha+\beta)^3-3\alpha\beta(\alpha+\beta)\\
  &=k^3-3(1)(k)\\
  &=k^3-3k
  \end{align*}
  So the product is
  \[
  (\alpha^2+\beta)(\beta^2+\alpha)=k^3-3k+2
  \]

  A quadratic with roots \(r_1\) and \(r_2\) is
  \[
  x^2-(r_1+r_2)x+r_1r_2=0
  \]
  Therefore the required equation is
  \[
  x^2-(k^2+k-2)x+(k^3-3k+2)=0
  \]

  So, in the form \(px^2+qx+r=0\),
  \[
  x^2-(k^2+k-2)x+(k^3-3k+2)=0
  \]
  with \(p=1\), \(q=-(k^2+k-2)\), \(r=k^3-3k+2\).
\end{enumerate}
\end{tcolorbox}


\newpage

% ---- Question 8 ----
\questionitem
% [Source: _QBT___Solns__Roots_of_Polynomials, Q8]

The roots of the equation
\[
x^3 + 5x^2 - 4x - 24 = 0
\]
are $\alpha$, $\beta$ and $\gamma$. \\

Without solving the equation,
\begin{enumerate}[label=\textbf{(\alph*)}]
  \item write down the value of each of
  \[
  \alpha + \beta + \gamma \qquad \alpha\beta + \beta\gamma + \gamma\alpha \qquad \alpha\beta\gamma
  \]
  \marks[-20pt]{1}
  \item Hence determine the value of
  \begin{enumerate}[label=(\roman*)]
    \item \[
    \frac{1}{\alpha\beta} + \frac{1}{\beta\gamma} + \frac{1}{\gamma\alpha}
    \]
    \marks[-20pt]{2}
    \item \[
    (\alpha + 2)(\beta + 2)(\gamma + 2)
    \]
    \marks[-20pt]{2}
    \item \[
    \alpha^2\beta^2 + \beta^2\gamma^2 + \gamma^2\alpha^2
    \]
    \marks[-20pt]{3}
  \end{enumerate}
\end{enumerate}

\vspace*{10pt}

\begin{tcolorbox}[colback=white, colframe=scarlet, title={\textbf{Solution}}, breakable, grow to left by=6mm, grow to right by=10mm]
\begin{enumerate}[label=\textbf{(\alph*)}]
  \item For a monic cubic
  \[
  x^3+bx^2+cx+d=0
  \]
  with roots $\alpha,\beta,\gamma$,
  \[
  \alpha+\beta+\gamma=-b,\qquad \alpha\beta+\beta\gamma+\gamma\alpha=c,\qquad \alpha\beta\gamma=-d
  \]
  Here
  \[
  x^3+5x^2-4x-24=0
  \]
  so
  \[
  \alpha+\beta+\gamma=-5,\qquad \alpha\beta+\beta\gamma+\gamma\alpha=-4,\qquad \alpha\beta\gamma=24
  \]
  Therefore the values are $-5$, $-4$ and $24$.

  \item Using
  \[
  \alpha+\beta+\gamma=-5,\qquad \alpha\beta+\beta\gamma+\gamma\alpha=-4,\qquad \alpha\beta\gamma=24
  \]

  \begin{enumerate}[label=(\roman*)]
    \item
    \begin{align*}
    \frac{1}{\alpha\beta}+\frac{1}{\beta\gamma}+\frac{1}{\gamma\alpha}
    &=\frac{\gamma}{\alpha\beta\gamma}+\frac{\alpha}{\alpha\beta\gamma}+\frac{\beta}{\alpha\beta\gamma} \\
    &=\frac{\alpha+\beta+\gamma}{\alpha\beta\gamma} \\
    &=\frac{-5}{24}
    \end{align*}
    Hence
    \[
    \frac{1}{\alpha\beta}+\frac{1}{\beta\gamma}+\frac{1}{\gamma\alpha}=-\frac{5}{24}
    \]

    \item
    \begin{align*}
    (\alpha+2)(\beta+2)(\gamma+2)
    &=(\alpha\beta+2\alpha+2\beta+4)(\gamma+2) \\
    &=\alpha\beta\gamma+2\alpha\beta+2\alpha\gamma+2\beta\gamma+4\alpha+4\beta+4\gamma+8 \\
    &=\alpha\beta\gamma+2(\alpha\beta+\beta\gamma+\gamma\alpha)+4(\alpha+\beta+\gamma)+8
    \end{align*}
    Substituting the values from part (a),
    \begin{align*}
    (\alpha+2)(\beta+2)(\gamma+2)
    &=24+2(-4)+4(-5)+8 \\
    &=24-8-20+8 \\
    &=4
    \end{align*}
    Therefore
    \[
    (\alpha+2)(\beta+2)(\gamma+2)=4
    \]

    \item Start with
    \[
    (\alpha\beta+\beta\gamma+\gamma\alpha)^2
    \]
    Expanding,
    \begin{align*}
    (\alpha\beta+\beta\gamma+\gamma\alpha)^2
    &=\alpha^2\beta^2+\beta^2\gamma^2+\gamma^2\alpha^2
    +2\alpha\beta\cdot\beta\gamma+2\beta\gamma\cdot\gamma\alpha+2\gamma\alpha\cdot\alpha\beta \\
    &=\alpha^2\beta^2+\beta^2\gamma^2+\gamma^2\alpha^2
    +2\alpha\beta^2\gamma+2\alpha\beta\gamma^2+2\alpha^2\beta\gamma \\
    &=\alpha^2\beta^2+\beta^2\gamma^2+\gamma^2\alpha^2
    +2\alpha\beta\gamma(\alpha+\beta+\gamma)
    \end{align*}
    So
    \begin{align*}
    \alpha^2\beta^2+\beta^2\gamma^2+\gamma^2\alpha^2
    &=(\alpha\beta+\beta\gamma+\gamma\alpha)^2-2\alpha\beta\gamma(\alpha+\beta+\gamma) \\
    &=(-4)^2-2(24)(-5) \\
    &=16+240 \\
    &=256
    \end{align*}
    Hence
    \[
    \alpha^2\beta^2+\beta^2\gamma^2+\gamma^2\alpha^2=256
    \]
  \end{enumerate}
\end{enumerate}
\end{tcolorbox}


\newpage

% ---- Question 9 ----
\questionitem
% [Source: _QBT___Solns__Roots_of_Polynomials, Q9]

The roots of the equation $3x^3 + px^2 - 12x - 8 = 0$ are $\alpha$, $\beta$ and $\gamma$. \\

\begin{enumerate}[label=\textbf{(\alph*)}]
  \item Given that $\alpha + \beta + \gamma = 2$, write down the value of $p$ \marks{1} \\
  \item Write down values for $\alpha\beta + \beta\gamma + \gamma\alpha$ and $\alpha\beta\gamma$ \marks{1} \\
  \item Hence find the value of $\left(1+\frac{1}{\alpha}\right)\left(1+\frac{1}{\beta}\right)\left(1+\frac{1}{\gamma}\right)$ \marks{3} \\
\end{enumerate}

\vspace*{10pt}

\begin{tcolorbox}[colback=white, colframe=scarlet, title={\textbf{Solution}}, breakable, grow to left by=6mm, grow to right by=10mm]
\begin{enumerate}[label=\textbf{(\alph*)}]
  \item For \(3x^3+px^2-12x-8=0\), the sum of the roots is
  \[
  \alpha+\beta+\gamma=-\frac{p}{3}
  \]
  Given \(\alpha+\beta+\gamma=2\),
  \[
  -\frac{p}{3}=2
  \]
  so
  \[
  p=-6
  \]
  Therefore, \(p=-6\).

  \item For \(3x^3+px^2-12x-8=0\),
  \[
  \alpha\beta+\beta\gamma+\gamma\alpha=\frac{-12}{3}=-4
  \]
  and
  \[
  \alpha\beta\gamma=-\frac{-8}{3}=\frac{8}{3}
  \]
  Therefore,
  \[
  \alpha\beta+\beta\gamma+\gamma\alpha=-4,\qquad \alpha\beta\gamma=\frac{8}{3}
  \]

  \item
  \begin{align*}
  \left(1+\frac{1}{\alpha}\right)\left(1+\frac{1}{\beta}\right)\left(1+\frac{1}{\gamma}\right)
  &=\frac{(\alpha+1)(\beta+1)(\gamma+1)}{\alpha\beta\gamma} \\
  &=\frac{\alpha\beta\gamma+\alpha\beta+\beta\gamma+\gamma\alpha+\alpha+\beta+\gamma+1}{\alpha\beta\gamma}
  \end{align*}
  Using the results from parts (a) and (b),
  \[
  \alpha+\beta+\gamma=2,\qquad
  \alpha\beta+\beta\gamma+\gamma\alpha=-4,\qquad
  \alpha\beta\gamma=\frac{8}{3}
  \]
  so
  \begin{align*}
  \left(1+\frac{1}{\alpha}\right)\left(1+\frac{1}{\beta}\right)\left(1+\frac{1}{\gamma}\right)
  &=\frac{\frac{8}{3}-4+2+1}{\frac{8}{3}} \\
  &=\frac{\frac{8}{3}-1}{\frac{8}{3}} \\
  &=\frac{\frac{5}{3}}{\frac{8}{3}} \\
  &=\frac{5}{8}
  \end{align*}
  Therefore,
  \[
  \left(1+\frac{1}{\alpha}\right)\left(1+\frac{1}{\beta}\right)\left(1+\frac{1}{\gamma}\right)=\frac{5}{8}
  \]
\end{enumerate}
\end{tcolorbox}


\newpage

% ---- Question 10 ----
\questionitem
% [Source: _QBT___Solns__Roots_of_Polynomials, Q10]

The roots of the equation $x^3 + 2x^2 - 5x - 3 = 0$ are $\alpha$, $\beta$ and $\gamma$. \\

\begin{enumerate}[label=\textbf{(\alph*)}]
  \item Find $\alpha^2\beta^2 + \beta^2\gamma^2 + \gamma^2\alpha^2$ \marks{4} \\
  \item Find an equation with integer coefficients whose roots are $\alpha^2$, $\beta^2$ and $\gamma^2$ \marks{4} \\
\end{enumerate}

\vspace*{10pt}

\begin{tcolorbox}[colback=white, colframe=scarlet, title={\textbf{Solution}}, breakable, grow to left by=6mm, grow to right by=10mm]
\begin{enumerate}[label=\textbf{(\alph*)}]
  \item Since $\alpha$, $\beta$ and $\gamma$ are roots of
\[
x^3+2x^2-5x-3=0
\]
the standard relations between roots and coefficients give
\[
\alpha+\beta+\gamma=-2,\qquad \alpha\beta+\beta\gamma+\gamma\alpha=-5,\qquad \alpha\beta\gamma=3
\]

Now expand
\begin{align*}
(\alpha\beta+\beta\gamma+\gamma\alpha)^2
&=\alpha^2\beta^2+\beta^2\gamma^2+\gamma^2\alpha^2
   +2(\alpha\beta)(\beta\gamma)+2(\beta\gamma)(\gamma\alpha)+2(\gamma\alpha)(\alpha\beta)\\
&=\alpha^2\beta^2+\beta^2\gamma^2+\gamma^2\alpha^2
   +2\alpha\beta^2\gamma+2\beta\gamma^2\alpha+2\gamma\alpha^2\beta\\
&=\alpha^2\beta^2+\beta^2\gamma^2+\gamma^2\alpha^2
   +2\alpha\beta\gamma(\alpha+\beta+\gamma)
\end{align*}

So
\begin{align*}
\alpha^2\beta^2+\beta^2\gamma^2+\gamma^2\alpha^2
&=(\alpha\beta+\beta\gamma+\gamma\alpha)^2-2\alpha\beta\gamma(\alpha+\beta+\gamma)\\
&=(-5)^2-2(3)(-2)\\
&=25+12\\
&=37
\end{align*}

Hence
\[
\alpha^2\beta^2+\beta^2\gamma^2+\gamma^2\alpha^2=37
\]

  \item Let the required equation have roots $\alpha^2$, $\beta^2$ and $\gamma^2$.

First find their sum:
\begin{align*}
\alpha^2+\beta^2+\gamma^2
&=(\alpha+\beta+\gamma)^2-2(\alpha\beta+\beta\gamma+\gamma\alpha)\\
&=(-2)^2-2(-5)\\
&=4+10\\
&=14
\end{align*}

From part (a),
\[
\alpha^2\beta^2+\beta^2\gamma^2+\gamma^2\alpha^2=37
\]

Also,
\[
\alpha^2\beta^2\gamma^2=(\alpha\beta\gamma)^2=3^2=9
\]

So for roots $\alpha^2,\beta^2,\gamma^2$, the monic cubic is
\[
x^3-(\text{sum of roots})x^2+(\text{sum of pairwise products})x-(\text{product of roots})=0
\]

Therefore
\[
x^3-14x^2+37x-9=0
\]

So the required equation is
\[
x^3-14x^2+37x-9=0
\]
\end{enumerate}
\end{tcolorbox}


\newpage

% ---- Question 11 ----
\questionitem
% [Source: _QBT___Solns__Roots_of_Polynomials, Q11]

The graph of $y=x^3+3x^2+px-20$, where $p$ is a constant, crosses the $x$-axis at three points whose $x$-coordinates are equally spaced. \\

Find the roots of
\[
x^3+3x^2+px-20=0
\] \marks[-20pt]{6}

\vspace*{10pt}

\begin{tcolorbox}[colback=white, colframe=scarlet, title={\textbf{Solution}}, breakable, grow to left by=6mm, grow to right by=10mm]
Since the three roots are equally spaced, let them be
\[
a-d,\quad a,\quad a+d
\]

For
\[
x^3+3x^2+px-20=0
\]
the sum of the roots is \(-3\), so
\[
(a-d)+a+(a+d)=-3
\]
\[
3a=-3
\]
\[
a=-1
\]

Also, the product of the roots is \(20\), because for a monic cubic \(x^3+\cdots-20\), the product of the roots is \(-(-20)=20\).

So
\[
(a-d)a(a+d)=20
\]

Now
\[
(a-d)(a+d)=a^2-d^2
\]
so
\[
a(a^2-d^2)=20
\]

Substitute \(a=-1\):
\[
(-1)\bigl(1-d^2\bigr)=20
\]
\[
-1+d^2=20
\]
\[
d^2=21
\]
\[
d=\sqrt{21}
\]

Therefore the roots are
\[
a-d=-1-\sqrt{21},\qquad a=-1,\qquad a+d=-1+\sqrt{21}
\]

Hence the roots of the equation are
\[
-1-\sqrt{21},\quad -1,\quad -1+\sqrt{21}
\]

(Checking: the pairwise sum is
\[
(-1-\sqrt{21})(-1)+(-1)(-1+\sqrt{21})+(-1-\sqrt{21})(-1+\sqrt{21})=-18
\]
so \(p=-18\), which is consistent.)
\end{tcolorbox}


\newpage

% ---- Question 12 ----
\questionitem
% [Source: _QBT___Solns__Roots_of_Polynomials, Q12]

The cubic equation $x^3-3x+1=0$ has roots $\alpha$, $\beta$ and $\gamma$. \\

\begin{enumerate}[label=\textbf{(\alph*)}]
  \item Find a cubic equation whose roots are $\dfrac{1}{\alpha-1}$, $\dfrac{1}{\beta-1}$, $\dfrac{1}{\gamma-1}$. \marks{3} \\
  \item Hence find the value of $\dfrac{1}{(\alpha-1)^2}+\dfrac{1}{(\beta-1)^2}+\dfrac{1}{(\gamma-1)^2}$. \marks{2} \\
  \item Find also the value of $\dfrac{1}{(\alpha-1)^5}+\dfrac{1}{(\beta-1)^5}+\dfrac{1}{(\gamma-1)^5}$. \marks{3} \\
\end{enumerate}

\vspace*{10pt}

\begin{tcolorbox}[colback=white, colframe=scarlet, title={\textbf{Solution}}, breakable, grow to left by=6mm, grow to right by=10mm]
\begin{enumerate}[label=\textbf{(\alph*)}]
  \item Let
\[
t=\frac{1}{x-1}
\]
so that
\[
x=1+\frac{1}{t}
\]

We substitute this into the given equation
\[
x^3-3x+1=0
\]

Then
\begin{align*}
\left(1+\frac{1}{t}\right)^3-3\left(1+\frac{1}{t}\right)+1&=0\\
\left(1+\frac{3}{t}+\frac{3}{t^2}+\frac{1}{t^3}\right)-3-\frac{3}{t}+1&=0\\
\frac{1}{t^3}+\frac{3}{t^2}-1&=0
\end{align*}

Multiply through by $t^3$:
\begin{align*}
1+3t-t^3&=0\\
t^3-3t-1&=0
\end{align*}

So the cubic equation whose roots are
\[
\frac{1}{\alpha-1},\quad \frac{1}{\beta-1},\quad \frac{1}{\gamma-1}
\]
is
\[
t^3-3t-1=0
\]

  \item Let the roots of
\[
t^3-3t-1=0
\]
be
\[
r_1=\frac{1}{\alpha-1},\quad r_2=\frac{1}{\beta-1},\quad r_3=\frac{1}{\gamma-1}
\]

For the cubic
\[
t^3+0t^2-3t-1=0
\]
we have
\[
r_1+r_2+r_3=0,\qquad r_1r_2+r_2r_3+r_3r_1=-3
\]

Now
\[
(r_1+r_2+r_3)^2=r_1^2+r_2^2+r_3^2+2(r_1r_2+r_2r_3+r_3r_1)
\]

So
\begin{align*}
r_1^2+r_2^2+r_3^2
&=(r_1+r_2+r_3)^2-2(r_1r_2+r_2r_3+r_3r_1)\\
&=0^2-2(-3)\\
&=6
\end{align*}

Hence
\[
\frac{1}{(\alpha-1)^2}+\frac{1}{(\beta-1)^2}+\frac{1}{(\gamma-1)^2}=6
\]

  \item Using the same roots $r_1,r_2,r_3$, each root satisfies
\[
r^3-3r-1=0
\]
so
\[
r^3=3r+1
\]

Multiply by $r^2$:
\[
r^5=r^2r^3=r^2(3r+1)=3r^3+r^2
\]

Now substitute again for $r^3$:
\begin{align*}
r^5&=3(3r+1)+r^2\\
&=9r+3+r^2
\end{align*}

Summing for $r_1,r_2,r_3$ gives
\begin{align*}
r_1^5+r_2^5+r_3^5
&=9(r_1+r_2+r_3)+3\cdot 3+(r_1^2+r_2^2+r_3^2)\\
&=9(0)+9+6\\
&=15
\end{align*}

Therefore
\[
\frac{1}{(\alpha-1)^5}+\frac{1}{(\beta-1)^5}+\frac{1}{(\gamma-1)^5}=15
\]
\end{enumerate}
\end{tcolorbox}


\newpage

% ---- Question 13 ----
\questionitem
% [Source: _QBT___Solns__Roots_of_Polynomials, Q13]

The quartic equation $x^4-2x^2+2x+1=0$ has roots $\alpha$, $\beta$, $\gamma$, $\delta$. \\

\begin{enumerate}[label=\textbf{(\alph*)}]
  \item Find a quartic equation whose roots are $\alpha^2$, $\beta^2$, $\gamma^2$, $\delta^2$ and state the value of $\alpha^2+\beta^2+\gamma^2+\delta^2$. \marks{5} \\
  \item Find the value of $\alpha^6+\beta^6+\gamma^6+\delta^6$. \marks{3} \\
  \item Find the value of $\alpha^8+\beta^8+\gamma^8+\delta^8$. \marks{2} \\
\end{enumerate}

\vspace*{10pt}

\begin{tcolorbox}[colback=white, colframe=scarlet, title={\textbf{Solution}}, breakable, grow to left by=6mm, grow to right by=10mm]
\begin{enumerate}[label=\textbf{(\alph*)}]
  \item Let
\[
f(x)=x^4-2x^2+2x+1
\]

If \(f(\alpha)=0\), then \(f(\alpha)f(-\alpha)=0\). So \(\alpha^2\) is a root of the equation obtained from \(f(x)f(-x)=0\) after writing everything in terms of \(x^2\). The same will then be true for \(\beta^2,\gamma^2,\delta^2\).

Now
\[
f(x)=(x^2-1)^2+2x,\qquad f(-x)=(x^2-1)^2-2x
\]

So
\begin{align*}
f(x)f(-x)
&=\bigl((x^2-1)^2+2x\bigr)\bigl((x^2-1)^2-2x\bigr)\\
&=(x^2-1)^4-4x^2
\end{align*}

Let \(y=x^2\). Then
\begin{align*}
(x^2-1)^4-4x^2
&=(y-1)^4-4y\\
&=y^4-4y^3+6y^2-4y+1-4y\\
&=y^4-4y^3+6y^2-8y+1
\end{align*}

Hence the required quartic equation is
\[
y^4-4y^3+6y^2-8y+1=0
\]

Its roots are \(\alpha^2,\beta^2,\gamma^2,\delta^2\).

If these roots are \(y_1,y_2,y_3,y_4\), then for a monic quartic,
\[
y_1+y_2+y_3+y_4=-(\text{coefficient of }y^3)=4
\]

Therefore
\[
\alpha^2+\beta^2+\gamma^2+\delta^2=4
\]

  \item Let the roots of
\[
h(y)=y^4-4y^3+6y^2-8y+1
\]
be \(y_1,y_2,y_3,y_4\), where
\[
y_1=\alpha^2,\quad y_2=\beta^2,\quad y_3=\gamma^2,\quad y_4=\delta^2
\]

Define
\[
q_r=y_1^r+y_2^r+y_3^r+y_4^r
\]

Then
\[
\alpha^6+\beta^6+\gamma^6+\delta^6=q_3
\]

For a monic quartic \(y^4+ay^3+by^2+cy+d\), the power sums satisfy
\begin{align*}
q_1+a&=0\\
q_2+aq_1+2b&=0\\
q_3+aq_2+bq_1+3c&=0
\end{align*}

Here \(a=-4,\ b=6,\ c=-8\), so
\begin{align*}
q_1-4&=0\\
q_2-4q_1+12&=0\\
q_3-4q_2+6q_1-24&=0
\end{align*}

From the first equation,
\[
q_1=4
\]

Then
\[
q_2-4(4)+12=0
\]
so
\[
q_2=4
\]

Now
\[
q_3-4(4)+6(4)-24=0
\]
which gives
\[
q_3=16
\]

Therefore
\[
\alpha^6+\beta^6+\gamma^6+\delta^6=16
\]

  \item Here
\[
\alpha^8+\beta^8+\gamma^8+\delta^8=q_4
\]

Using the next power-sum relation for the same quartic,
\[
q_4-4q_3+6q_2-8q_1+4=0
\]

Substitute \(q_1=4,\ q_2=4,\ q_3=16\):
\begin{align*}
q_4-4(16)+6(4)-8(4)+4&=0\\
q_4-64+24-32+4&=0\\
q_4-68&=0
\end{align*}

So
\[
q_4=68
\]

Hence
\[
\alpha^8+\beta^8+\gamma^8+\delta^8=68
\]
\end{enumerate}
\end{tcolorbox}


\newpage

% ---- Question 14 ----
\questionitem
% [Source: _QBT___Solns__Roots_of_Polynomials, Q14]

The equation $z^4-2z^3-z^2-2z+1=0$ has roots $\alpha$, $\beta$, $\gamma$ and $\delta$. \\

\begin{enumerate}[label=\textbf{(\alph*)}]
  \item Show that a quartic equation whose roots are $\alpha+\dfrac{1}{\alpha}$, $\beta+\dfrac{1}{\beta}$, $\gamma+\dfrac{1}{\gamma}$ and $\delta+\dfrac{1}{\delta}$ is
  \[
  w^4-4w^3-2w^2+12w+9=0
  \]
  \marks[-20pt]{3}
  \item Hence determine the exact roots of the equation $z^4-2z^3-z^2-2z+1=0$. \marks{3} \\
\end{enumerate}

\vspace*{10pt}

\begin{tcolorbox}[colback=white, colframe=scarlet, title={\textbf{Solution}}, breakable, grow to left by=6mm, grow to right by=10mm]
\begin{enumerate}[label=\textbf{(\alph*)}]
  \item Let
\[
f(z)=z^4-2z^3-z^2-2z+1
\]
and let $z$ be any root of $f(z)=0$.

Since the constant term is $1$, no root is $0$, so we may divide by $z^2$:
\begin{align*}
z^2-2z-1-\frac{2}{z}+\frac{1}{z^2}&=0 \\
\left(z^2+\frac{1}{z^2}\right)-2\left(z+\frac{1}{z}\right)-1&=0
\end{align*}

Now let
\[
w=z+\frac{1}{z}
\]
Then
\[
z^2+\frac{1}{z^2}=\left(z+\frac{1}{z}\right)^2-2=w^2-2
\]
so
\begin{align*}
(w^2-2)-2w-1&=0 \\
w^2-2w-3&=0 \\
(w-3)(w+1)&=0
\end{align*}

So the possible values of $w$ are $3$ and $-1$.

Also, the polynomial is unchanged when its coefficients are reversed, so if $z$ is a root then $\frac1z$ is also a root. Therefore the four roots $\alpha,\beta,\gamma,\delta$ occur in reciprocal pairs, and the four values
\[
\alpha+\frac1\alpha,\quad \beta+\frac1\beta,\quad \gamma+\frac1\gamma,\quad \delta+\frac1\delta
\]
are
\[
3,3,-1,-1
\]

Hence the required quartic is
\begin{align*}
(w-3)^2(w+1)^2&=(w^2-2w-3)^2 \\
&=w^4-4w^3-2w^2+12w+9
\end{align*}

So the quartic equation is
\[
w^4-4w^3-2w^2+12w+9=0
\]

  \item From part (a), the corresponding values of
\[
w=z+\frac1z
\]
are $3$ and $-1$.

If $w=3$, then
\begin{align*}
z+\frac1z&=3 \\
z^2-3z+1&=0
\end{align*}
So
\[
z=\frac{3\pm\sqrt5}{2}
\]

If $w=-1$, then
\begin{align*}
z+\frac1z&=-1 \\
z^2+z+1&=0
\end{align*}
So
\[
z=\frac{-1\pm \mathrm{i}\sqrt3}{2}
\]

Hence the exact roots of
\[
z^4-2z^3-z^2-2z+1=0
\]
are
\[
\frac{3+\sqrt5}{2},\quad \frac{3-\sqrt5}{2},\quad \frac{-1+\mathrm{i}\sqrt3}{2},\quad \frac{-1-\mathrm{i}\sqrt3}{2}
\]
\end{enumerate}
\end{tcolorbox}


\newpage

% ---- Question 15 ----
\questionitem
% [Source: _QBT___Solns__Roots_of_Polynomials, Q15]

The quartic equation $x^4-2x^3+3x^2-4x+1=0$ has roots $\alpha$, $\beta$, $\gamma$, $\delta$. \\

\begin{enumerate}[label=\textbf{(\alph*)}]
  \item Find the value of $\dfrac{1}{\alpha}+\dfrac{1}{\beta}+\dfrac{1}{\gamma}+\dfrac{1}{\delta}$. \marks{2} \\
  \item Find the value of $\dfrac{1}{\alpha^2}+\dfrac{1}{\beta^2}+\dfrac{1}{\gamma^2}+\dfrac{1}{\delta^2}$. \marks{2} \\
  \item Find the value of
  \[
  \left(\dfrac{1}{\alpha}+\dfrac{1}{\beta}+\dfrac{1}{\gamma}\right)^2+\left(\dfrac{1}{\beta}+\dfrac{1}{\gamma}+\dfrac{1}{\delta}\right)^2+\left(\dfrac{1}{\gamma}+\dfrac{1}{\delta}+\dfrac{1}{\alpha}\right)^2+\left(\dfrac{1}{\delta}+\dfrac{1}{\alpha}+\dfrac{1}{\beta}\right)^2
  \]
  \marks[-20pt]{5}
\end{enumerate}

\vspace*{10pt}

\begin{tcolorbox}[colback=white, colframe=scarlet, title={\textbf{Solution}}, breakable, grow to left by=6mm, grow to right by=10mm]
\begin{enumerate}[label=\textbf{(\alph*)}]
  \item For the monic quartic
  \[
  x^4-2x^3+3x^2-4x+1=0
  \]
  with roots \(\alpha,\beta,\gamma,\delta\), the standard root relations give
  \[
  \alpha+\beta+\gamma+\delta=2
  \]
  \[
  \alpha\beta+\alpha\gamma+\alpha\delta+\beta\gamma+\beta\delta+\gamma\delta=3
  \]
  \[
  \alpha\beta\gamma+\alpha\beta\delta+\alpha\gamma\delta+\beta\gamma\delta=4
  \]
  \[
  \alpha\beta\gamma\delta=1
  \]

  Hence
  \begin{align*}
  \frac{1}{\alpha}+\frac{1}{\beta}+\frac{1}{\gamma}+\frac{1}{\delta}
  &=\frac{\beta\gamma\delta+\alpha\gamma\delta+\alpha\beta\delta+\alpha\beta\gamma}{\alpha\beta\gamma\delta}\\
  &=\frac{4}{1}\\
  &=4
  \end{align*}

  So
  \[
  \frac{1}{\alpha}+\frac{1}{\beta}+\frac{1}{\gamma}+\frac{1}{\delta}=4
  \]

  \item Let
  \[
  r_1=\frac{1}{\alpha},\quad r_2=\frac{1}{\beta},\quad r_3=\frac{1}{\gamma},\quad r_4=\frac{1}{\delta}
  \]

  Then from part (a),
  \[
  r_1+r_2+r_3+r_4=4
  \]

  Also
  \begin{align*}
  r_1r_2+r_1r_3+r_1r_4+r_2r_3+r_2r_4+r_3r_4
  &=\frac{1}{\alpha\beta}+\frac{1}{\alpha\gamma}+\frac{1}{\alpha\delta}+\frac{1}{\beta\gamma}+\frac{1}{\beta\delta}+\frac{1}{\gamma\delta}\\
  &=\frac{\gamma\delta+\beta\delta+\beta\gamma+\alpha\delta+\alpha\gamma+\alpha\beta}{\alpha\beta\gamma\delta}\\
  &=\frac{\alpha\beta+\alpha\gamma+\alpha\delta+\beta\gamma+\beta\delta+\gamma\delta}{\alpha\beta\gamma\delta}\\
  &=\frac{3}{1}\\
  &=3
  \end{align*}

  Now use
  \[
  (r_1+r_2+r_3+r_4)^2=r_1^2+r_2^2+r_3^2+r_4^2+2(r_1r_2+r_1r_3+r_1r_4+r_2r_3+r_2r_4+r_3r_4)
  \]

  So
  \begin{align*}
  \frac{1}{\alpha^2}+\frac{1}{\beta^2}+\frac{1}{\gamma^2}+\frac{1}{\delta^2}
  &=4^2-2(3)\\
  &=16-6\\
  &=10
  \end{align*}

  So
  \[
  \frac{1}{\alpha^2}+\frac{1}{\beta^2}+\frac{1}{\gamma^2}+\frac{1}{\delta^2}=10
  \]

  \item Using the same notation, let
  \[
  S=r_1+r_2+r_3+r_4=4
  \]

  Then each bracket is the sum of all four reciprocals except one:
  \[
  \frac{1}{\alpha}+\frac{1}{\beta}+\frac{1}{\gamma}=S-r_4
  \]
  \[
  \frac{1}{\beta}+\frac{1}{\gamma}+\frac{1}{\delta}=S-r_1
  \]
  \[
  \frac{1}{\gamma}+\frac{1}{\delta}+\frac{1}{\alpha}=S-r_2
  \]
  \[
  \frac{1}{\delta}+\frac{1}{\alpha}+\frac{1}{\beta}=S-r_3
  \]

  So the required expression is
  \[
  (S-r_1)^2+(S-r_2)^2+(S-r_3)^2+(S-r_4)^2
  \]

  Expanding,
  \begin{align*}
  \sum_{i=1}^4 (S-r_i)^2
  &=\sum_{i=1}^4 \left(S^2-2Sr_i+r_i^2\right)\\
  &=4S^2-2S\sum_{i=1}^4 r_i+\sum_{i=1}^4 r_i^2
  \end{align*}

  From parts (a) and (b),
  \[
  \sum_{i=1}^4 r_i=S=4
  \qquad\text{and}\qquad
  \sum_{i=1}^4 r_i^2=10
  \]

  Therefore
  \begin{align*}
  \sum_{i=1}^4 (S-r_i)^2
  &=4(4^2)-2(4)(4)+10\\
  &=64-32+10\\
  &=42
  \end{align*}

  So the value of the expression is
  \[
  42
  \]
\end{enumerate}
\end{tcolorbox}


\end{enumerate}
\end{document}
